\section{Experimental Method}
\label{sec:experimental-method}
\subsection{Workload and modeled platform}
FanOutQA questions, human decompositions and evidence~\cite{zhu2024fanoutqa}
deterministically produce 160 complete tasks and 1,304 calls. A pinned
Qwen3-8B tokenizer~\cite{qwen2025technical} preserves prompt identities;
durations and outputs are prescribed. Context reaches 86,214 tokens without
truncation or padding; the 128K bucket is empty. The 131,072-token proposed-YaRN
capacity validates neither extended-context numerics nor agent answer quality.

\begin{table}[t]
\centering
\caption{Declared simulation inputs; no row is an A100 measurement.}
\label{tab:configuration}
\setlength{\tabcolsep}{3pt}
\begin{tabular}{@{}ll@{}}
\toprule
Parameter & Value \\
\midrule
Model / DP / TP / PP & Qwen3-8B / 8 / 1 / 1 \\
HBM / weights per GPU & 40 GiB / 16,381,470,720 B \\
Fixed runtime reserve & 2 GiB per GPU \\
Registered HBM capacity $M$ & 40 GiB per target \\
DRAM Store & $10^{12}$ B, one colocated node \\
Effective path service rate & 25 GB/s (decimal) \\
Transfer setup / reroute & 250 $\mu$s / 1 ms \\
Protected prompt chunk & 4,096 tokens \\
Physical allocation unit & 64 tokens \\
Workflow concurrency $C$ & 16 (also 8, 32) \\
Initial / maximum $\alpha$ & 0.20 / 0.50 \\
AIMD $\gamma$ / $\delta$ / $K$ & 0.50 / 0.05 / 4 \\
Control interval & 1 s \\
Detection delay / repair & 0 / 120 s \\
\bottomrule
\end{tabular}
\end{table}

Figure~\ref{fig:workload} describes the workload. The arms are unprotected
recomputation, AegisKV DRAM-only and full HBM+DRAM. DRAM-only retains
SlackLandlord; it represents neither vanilla Mooncake nor GhostServe.
Every enabled tier sees every sealed prompt candidate. Salts, routing,
service models, workload and physical capacities are common.

\begin{figure*}[tp]
\centering\includegraphics[width=\textwidth]{figures/aligned-workload.pdf}
\caption{The frozen workload. Prompt sizes
and calls per workflow are descriptive ECDFs. Context coverage uses exact
prompt tokens plus prescribed outputs; the 128K bucket contains no requests.}
\Description{Workload size distributions and context bucket counts.}
\label{fig:workload}
\end{figure*}

\subsection{Pairing and interventions}
Closed-loop admission keeps at most $C$ workflows active and admits the next
one when a workflow completes. Five SHA-256-based admission permutations
(seeds 0--4) are paired across treatments and no-failure controls. Routing
uses a fixed seed of zero. WCT begins at a workflow's first request release,
excluding its wait for a workflow-admission slot. Complete-workload makespan
also exposes that waiting through the final completion time.

Before measured runs, an unprotected reference selects eligible 96K-bucket
requests near 50\% workflow progress (one fault) or one-third/two-thirds
(two faults). Sequential plans incorporate previously selected faults;
selection never conditions on protection hits or benefit. Each fault targets
the GPU naturally serving the selected request when it crosses 60\% of
prompt-plus-output progress. The event interpolates within the crossing
scheduler step, but only earlier completed forwards expose KV. Physical
routes, fault times and collateral victims can differ across policies.
This estimates a matched logical-request intervention, not an identical
wall-clock shock. Repair takes 120 modeled seconds.

The base plan is reused for bandwidth and service-model sensitivities.
Concurrency and repeated-fault settings derive their own reference plans.
The complete prespecified matrix contains 170 measured executions:
60 base runs (six arms, F0/F1, five seeds), 40 concurrency runs, 40 service-model
runs, 20 bandwidth runs and 10 two-fault runs. The latter reuse the base F0
controls for subtraction. Selection runs are excluded from the result count.

\subsection{Ablations, timing and outcomes}
The three ablations remove the Slack signal, replace retention with LRU, or
disable AIMD while keeping the initial HBM quota. They retain request ownership,
copy/commit semantics, full candidate scope and completion cleanup. LRU uses
ordinary retention in place of Algorithm 2; it is not described as native
SlackLandlord or credited with that algorithm's suffix-victim guarantee.

Baseline batches compose participant service ideally. Sensitivities
interpolate toward serialization ($\beta=0.5,1$), preserving scheduler and
token sequences, or reduce path bandwidth from 25 to 5 GB/s. These vary
uncalibrated assumptions; exact timing equations accompany the artifact.

For seed $s$, treatment $a$, workflow $i$ and condition $F$, let
$W_{s,a,i}^{F}$ be WCT. We report healthy cost, raw F1 completion and the
paired difference in failure penalties:
\begin{equation}
\begin{split}
\Gamma_{s,i}={}&(W_{s,\mathrm{full},i}^{1}-W_{s,\mathrm{full},i}^{0})\\
 &-(W_{s,\mathrm{dram},i}^{1}-W_{s,\mathrm{dram},i}^{0}).
\end{split}
\end{equation}
Negative $\Gamma$ favors full. For an ablation, the contrast is its penalty
minus full, so positive values favor full. We first aggregate all 160
workflows within a seed. Focus-workflow contrasts use the prespecified
failure requests' workflows; collateral cohorts remain descriptive
post-treatment populations.

Two-sided 95\% Student-$t$ intervals use the five seed estimates ($df=4$).
The independent unit is an admission permutation of this fixed workload,
not each request or workflow. Intervals are exploratory and unadjusted for
multiple comparisons. All five observations remain visible. ECDFs describe
the finite workload; nearest-rank percentiles use $\lceil qn\rceil$. Restore breakdowns conserve
prompt tokens across local primary, copied HBM, DRAM and recomputation.
Physical pilot observations are analyzed separately.
